% ============================================================
\section{Comparison to Prior Protocols}
\label{sec:comparison}

M\textsuperscript{2}'s contribution is best assessed against the design space of prior on-chain protocols that have engaged, in some form, with the question of contractually enforceable claims on a treasury or peg target. Table~\ref{tab:protocol-comparison} reports the comparison along eight dimensions chosen to expose where M\textsuperscript{2}'s design diverges. The five DeFi comparators---Olympus DAO \citep{OlympusDAO2021}, Reflexer RAI \citep{Reflexer2021}, Frax \citep{Frax2021}, Liquity \citep{Liquity2021}, and MakerDAO Dai \citep{Maker2019}---span the canonical lineage of treasury-backed and CDP-style protocols. The TradFi closed-end-fund row provides the cross-domain anchor.

\begin{table}[t]
\centering
\scriptsize
\caption{Comparison matrix across eight design dimensions. ``Fixed supply'' refers to whether the protocol can mint new units of the primary asset post-genesis. ``On-chain redemption'' is a contractually enforceable exit at a protocol-defined price (versus secondary-market only). ``Floor invariant'' indicates whether a per-unit reserve claim is non-decreasing by contract. ``Operator dep.'' indicates whether the protocol's growth (not its safety) depends on a centralized off-chain actor.}
\label{tab:protocol-comparison}
\renewcommand{\arraystretch}{1.18}
\begin{tabular}{p{1.5cm} p{1.25cm} p{1.55cm} p{1.4cm} p{1.5cm} p{1.4cm} p{1.0cm} p{1.85cm}}
\toprule
\textbf{Protocol} & \textbf{Fixed supply?} & \textbf{On-chain redemption?} & \textbf{Treasury composition} & \textbf{Floor invariant?} & \textbf{Gov.} & \textbf{Op. dep.} & \textbf{Primary risk} \\
\midrule
M\textsuperscript{2} & Yes; supply only decreases & Yes; at $\Floor = \Tres/\Sup$ & Single exog.\ stable; one-way in & Yes; non-decreasing by Thm~\ref{thm:floor-monotonicity} & None & Growth only & Smart-contract bug; stable issuer \\
OHM (OlympusDAO) & No; staking rebases up & No (governance vote-gated burn proposals only) & Multi-asset; rebalance by governance & No (``RFV per token'' is rhetorical) & Active gov. & Low (community treasury) & Confidence collapse; gov.\ capture \\
RAI (Reflexer) & No; PID mint/burn & Yes; against ETH at redemption price & ETH collateral via Safe vaults & No (redemption price mean-reverts) & Minimal & None & Oracle manip.; vault liquidation cascade \\
Frax & No; algorithmic + collateralized & Yes; against AMO-managed reserves & Multi-stable + algorithmic & No (peg target, not floor) & Active gov. & Mod. & Peg loss; AMO mispricing \\
Liquity (LUSD) & No; mintable against ETH & Yes; against ETH at \$1 + 0.5\% fee & ETH-only via troves & No (peg target, not floor) & None & None & ETH price crash + recovery-mode cascade \\
MakerDAO (Dai) & No; CDP-style mint & Yes; CDP unwind only & Multi-collateral & Limited (per-vault) & Active gov. & Low & Black-swan oracle deviation; gov.\ capture \\
\midrule
TradFi CEF (closed-end fund) & Effectively yes & No (secondary market only) & Multi-asset, manager-discretion & No (NAV fluctuates; persistent discount) & Board / sh. & Yes (manager) & Discount-to-NAV persistence; mgmt.\ extraction \\
\bottomrule
\end{tabular}
\end{table}

\paragraph{What the table shows.} Three dimensions sharply discriminate M\textsuperscript{2} from its closest comparators. First, \emph{floor invariance}: M\textsuperscript{2} is the only entry in Table~\ref{tab:protocol-comparison} for which a contractually enforced, non-decreasing per-unit reserve claim is the central design property. OHM advertises ``risk-free value per token'' but does not enforce redemption against it; RAI's redemption price is mean-reverting under PID control, not monotone; Frax targets a peg; Liquity targets a \$1 peg via redemption at par rather than a monotone floor; Dai has per-vault redemption guarantees but no aggregate-level floor; classical CEFs publish a NAV that has no contractual redemption. Second, \emph{supply discipline}: only M\textsuperscript{2} guarantees that supply can only decrease post-genesis. OHM's staking rebases dilute holders; Frax and Maker can mint to absorb collateral; RAI mints and burns under PID control; Liquity mints LUSD against new troves. Holders of an M\textsuperscript{2} token are guaranteed, under the protocol-defined operation set, that their fractional claim on the treasury cannot be diluted by post-genesis issuance. Third, \emph{governance surface}: M\textsuperscript{2} and Liquity are the only entries with no governance mechanism whatsoever. Liquity's immutability is a well-known precedent and a direct ancestor of the design stance taken here; M\textsuperscript{2} differs from Liquity in that its primary asset is not a stablecoin and its reserve is denominated in an external stable rather than an internal CDP-collateral pool.

\paragraph{What the table does not capture.} Several dimensions of comparison are inherently qualitative and resist the matrix format. The maturity of the deployed codebase: OHM, Frax, Liquity, and Dai have multi-year mainnet histories with large TVLs; M\textsuperscript{2} is pre-deployment at the time of writing. The depth of the user community and ecosystem integrations: the older protocols are deeply embedded in DeFi compositional stacks; M\textsuperscript{2} is not. The empirical reliability of the floor under stress: M\textsuperscript{2} has not yet faced a real-world adversary. We acknowledge these as legitimate concerns and address them programmatically in Section~\ref{sec:limitations} (formal-verification path) and in the empirical follow-up paper sketched in Section~\ref{sec:conclusion}.

\paragraph{The cross-domain analog.} The closed-end-fund row is the most informative TradFi comparator because of the structural similarity in capital structure (a pool of assets held in custody, claimed pro rata by a fixed number of units) and the structural difference in redemption mechanism. The persistent CEF discount-to-NAV documented by \citet{LeeShleiferThaler1991} and \citet{Pontiff1996} arises precisely because TradFi CEFs lack a contractual NAV-redemption interface; investors who want NAV must sell to another investor, who has no obligation to pay it. The cross-domain claim we make is that M\textsuperscript{2} is the implementation of NAV-redemption that the CEF literature has identified as the discount-collapsing primitive but that the TradFi legal substrate cannot enforce at the contract level (a CEF could in principle adopt mandatory NAV-redemption but then becomes an open-end fund subject to a different regulatory regime). The smart-contract substrate dissolves the structural obstacle, and M\textsuperscript{2} adds the further property of NAV monotonicity, which has no TradFi analog at all.
